%=========================
% Sections/09-DanhGiaHeThong.tex
% Chương 9 – Đánh giá hệ thống
%=========================

\section{Đánh giá hệ thống}

Chương này trình bày cách đánh giá hệ thống sau khi cài đặt. Việc đánh giá cần dựa trên mục tiêu đề tài, yêu cầu chức năng, yêu cầu phi chức năng và các tiêu chí chất lượng đã xác định.

\subsection*{Mục tiêu đánh giá}

Nêu rõ mục tiêu của quá trình đánh giá, chẳng hạn kiểm tra mức độ hoàn thành chức năng, độ ổn định, hiệu năng, khả năng sử dụng hoặc mức độ đáp ứng yêu cầu nghiệp vụ.

\subsection*{Kịch bản đánh giá}

Mô tả các kịch bản, dữ liệu thử nghiệm, môi trường đánh giá và cách thu thập kết quả. Các kịch bản cần phản ánh các tình huống sử dụng quan trọng của hệ thống.

\begin{longtable}{@{}p{0.14\linewidth}p{0.34\linewidth}p{0.24\linewidth}p{0.18\linewidth}@{}}
\caption{Danh sách kịch bản đánh giá}\label{tab:evaluation-scenarios}\\
\toprule
\textbf{Mã} & \textbf{Kịch bản} & \textbf{Dữ liệu/Kiểm kiện} & \textbf{Kết quả kỳ vọng} \\
\midrule
\endfirsthead
\toprule
\textbf{Mã} & \textbf{Kịch bản} & \textbf{Dữ liệu/Kiểm kiện} & \textbf{Kết quả kỳ vọng} \\
\midrule
\endhead
EV-01 & \todoitem{Bổ sung kịch bản đánh giá} & \todoitem{Bổ sung} & \todoitem{Bổ sung} \\
EV-02 & \todoitem{Bổ sung kịch bản đánh giá} & \todoitem{Bổ sung} & \todoitem{Bổ sung} \\
\bottomrule
\end{longtable}

\subsection*{Kết quả đánh giá}

Trình bày kết quả thu được theo từng tiêu chí. Có thể sử dụng bảng, biểu đồ hoặc hình minh họa để thể hiện kết quả rõ ràng.

\begin{table}[H]
\centering
\caption{Mẫu tổng hợp kết quả đánh giá}
\label{tab:evaluation-results}
\begin{tabularx}{\linewidth}{@{}p{0.24\linewidth}p{0.24\linewidth}X@{}}
\toprule
\textbf{Tiêu chí} & \textbf{Kết quả} & \textbf{Nhận xét} \\
\midrule
Chức năng & \todoitem{Bổ sung} & \todoitem{Bổ sung nhận xét} \\
Hiệu năng & \todoitem{Bổ sung} & \todoitem{Bổ sung nhận xét} \\
Khả dụng & \todoitem{Bổ sung} & \todoitem{Bổ sung nhận xét} \\
\bottomrule
\end{tabularx}
\end{table}

\subsection*{Thảo luận}

Phân tích ý nghĩa của kết quả đánh giá, nguyên nhân của các điểm mạnh/yếu và mức độ đáp ứng mục tiêu ban đầu. Phần này nên liên kết lại với các yêu cầu ở Chương 5 và thiết kế ở Chương 7.

\subsection*{Kết chương}

Chương này đã trình bày mục tiêu, kịch bản và kết quả đánh giá hệ thống. Kết quả đánh giá là cơ sở để tổng kết kết quả đạt được, hạn chế và hướng phát triển ở Chương 10.
